\documentclass[11pt]{article}
\usepackage{fontspec,amsmath,amssymb,geometry}
\setmainfont{Times New Roman}\geometry{margin=0.85in}
\newcommand{\solutionquestion}[1]{\par\leftskip=0pt\section*{Q#1}\leftskip=1em}
\newcommand{\solutionpart}[1]{\par\smallskip\leftskip=1em\noindent\hangindent=2em\hangafter=1\makebox[1.5em][l]{\textbf{(#1)}}\hspace{0.5em}\ignorespaces}
\newcommand{\solutioncontinuation}{\par\leftskip=3em\noindent\ignorespaces}
\title{Solving Linear Equations}\author{T.T.}\date{}
\begin{document}\maketitle

\solutionquestion{1}
\solutionpart{i}The first pivot is $2$. To eliminate $x$ from the second equation, subtract five times the first equation from the second:
\[
\begin{aligned}
2x+3y&=1,\\
(10x+9y)-5(2x+3y)&=11-5.
\end{aligned}
\]
Therefore the upper triangular system is
\[
2x+3y=1,\qquad -6y=6.
\]
The two pivots are $2$ and $-6$.

\solutionpart{ii}Back substitution gives $y=-1$. The first equation then gives
\[
2x+3(-1)=1,
\]
so $x=2$. In column form,
\[
2\begin{bmatrix}2\\10\end{bmatrix}
-\begin{bmatrix}3\\9\end{bmatrix}
=\begin{bmatrix}1\\11\end{bmatrix},
\]
which verifies the solution.

\solutionpart{iii}The coefficient matrix and the elimination multiplier are unchanged. For the new right side, the second transformed entry is
\[
44-5(4)=24.
\]
Thus $-6y=24$, so $y=-4$. Substitution into $2x+3y=4$ gives
\[
2x-12=4,
\]
and hence $x=8$. The new solution is $(x,y)=(8,-4)$.

\solutionquestion{2}
\solutionpart{i}Since $x_1$ and $x_2$ are solutions,
\[
Ax_1=b,\qquad Ax_2=b.
\]
For any $t\in\mathbb R$, linearity gives
\[
\begin{aligned}
A\bigl((1-t)x_1+tx_2\bigr)
&=(1-t)Ax_1+tAx_2\\
&=(1-t)b+tb=b.
\end{aligned}
\]
Therefore $(1-t)x_1+tx_2$ is also a solution.

\solutionpart{ii}Because $x_1\ne x_2$, the vectors $(1-t)x_1+tx_2$ are distinct for distinct values of $t$. Indeed,
\[
(1-t_1)x_1+t_1x_2=(1-t_2)x_1+t_2x_2
\]
would imply
\[
(t_1-t_2)(x_2-x_1)=0,
\]
and therefore $t_1=t_2$. Hence two distinct solutions generate an entire line of solutions. A linear system can have no solution, one solution, or infinitely many solutions, but it cannot have exactly two.

\solutionquestion{3}
\solutionpart{i}The multiplier for the entry $a_{21}=4$ is $4$, so
\[
E_{21}=\begin{bmatrix}
1&0&0\\
-4&1&0\\
0&0&1
\end{bmatrix}.
\]
\solutioncontinuation{} The multiplier for $a_{31}=-2$ is $-2$. Subtracting $-2$ times row 1 is the same as adding twice row 1, so
\[
E_{31}=\begin{bmatrix}
1&0&0\\
0&1&0\\
2&0&1
\end{bmatrix}.
\]
\solutioncontinuation{} After these two steps,
\[
E_{31}E_{21}A
=\begin{bmatrix}
1&1&0\\
0&2&1\\
0&4&0
\end{bmatrix}.
\]
\solutioncontinuation{} The final multiplier is $4/2=2$, and therefore
\[
E_{32}=\begin{bmatrix}
1&0&0\\
0&1&0\\
0&-2&1
\end{bmatrix}.
\]

\solutionpart{ii}Multiplying the elimination matrices in the order in which they act gives
\[
M=E_{32}E_{31}E_{21}
=\begin{bmatrix}
1&0&0\\
-4&1&0\\
10&-2&1
\end{bmatrix}.
\]
Then
\[
MA=
\begin{bmatrix}
1&0&0\\
-4&1&0\\
10&-2&1
\end{bmatrix}
\begin{bmatrix}
1&1&0\\
4&6&1\\
-2&2&0
\end{bmatrix}
=\begin{bmatrix}
1&1&0\\
0&2&1\\
0&0&-2
\end{bmatrix}
=U.
\]

\solutionpart{iii}The transformed right side is
\[
Mb=
\begin{bmatrix}
1&0&0\\
-4&1&0\\
10&-2&1
\end{bmatrix}
\begin{bmatrix}1\\0\\0\end{bmatrix}
=\begin{bmatrix}1\\-4\\10\end{bmatrix}.
\]
Thus the triangular system is
\[
x+y=1,\qquad 2y+z=-4,\qquad -2z=10.
\]
Back substitution gives
\[
z=-5,\qquad y=\frac12,\qquad x=\frac12.
\]
Therefore
\[
x=\begin{bmatrix}\tfrac12\\[2pt]\tfrac12\\[2pt]-5\end{bmatrix}.
\]

\solutionquestion{4}
\solutionpart{i}Direct multiplication gives
\[
AB=\begin{bmatrix}7&0\\0&0\end{bmatrix},
\qquad
BA=\begin{bmatrix}1&2\\3&6\end{bmatrix}.
\]
Since $AB\ne BA$, these matrices do not commute.

\solutionpart{ii}Distributivity must be used without changing the order of the factors:
\[
{(A+B)}^{2}=(A+B)(A+B)=A^2+AB+BA+B^2.
\]

\solutionpart{iii}First,
\[
A+B=\begin{bmatrix}2&2\\3&0\end{bmatrix},
\qquad
{(A+B)}^{2}=\begin{bmatrix}10&4\\6&6\end{bmatrix}.
\]
Also $A^2=A$, $B^2=B$, and therefore
\[
A^2+2AB+B^2
=\begin{bmatrix}16&2\\3&0\end{bmatrix}.
\]
The two results are different. In general,
\[
A^2+AB+BA+B^2=A^2+2AB+B^2
\]
if and only if $BA=AB$. Thus the scalar-style identity is valid precisely when $A$ and $B$ commute.

\solutionquestion{5}
\solutionpart{i}Write the rows of $A$ as $r_1^{T},r_2^{T},r_3^{T}$, where
\[
r_3^{T}=r_1^{T}+r_2^{T}.
\]
If $Ax=e_3$, where $e_3={(0,0,1)}^{T}$, then its first two equations require
\[
r_1^{T}x=0,\qquad r_2^{T}x=0,
\]
while its third equation requires $r_3^{T}x=1$. However,
\[
r_3^{T}x=(r_1^{T}+r_2^{T})x=r_1^{T}x+r_2^{T}x=0,
\]
which is a contradiction. Thus $Ax=e_3$ has no solution, so $A$ cannot be invertible.

\solutioncontinuation{} For a general right side $b$, the same row relation requires
\[
b_3=b_1+b_2.
\]
Indeed, the elimination step $R_3\leftarrow R_3-R_1-R_2$ produces
\[
0=b_3-b_1-b_2.
\]
If this condition fails, the system is inconsistent. If it holds, the transformed third equation is $0=0$ and is redundant. The condition is necessary; further row relations could impose additional conditions.

\solutionpart{ii}Write the columns of $A$ as $a_1,a_2,a_3$, with $a_3=a_1+a_2$. For
\[
x=\begin{bmatrix}-1\\-1\\1\end{bmatrix},
\]
we have
\[
Ax=-a_1-a_2+a_3=0.
\]
This is a nonzero solution of $Ax=0$. If $A$ were invertible, multiplication by $A^{-1}$ would imply $x=A^{-1}0=0$, a contradiction. Therefore $A$ is not invertible.

\solutionquestion{6}
Start with the augmented matrix
\[
[A\ I]=
\left[
\begin{array}{ccc@{\;\vline\;}ccc}
2&1&0&1&0&0\\
1&2&1&0&1&0\\
0&1&2&0&0&1
\end{array}
\right].
\]
Apply $R_2\leftarrow2R_2$ followed by $R_2\leftarrow R_2-R_1$. Then apply $R_3\leftarrow3R_3$ followed by $R_3\leftarrow R_3-R_2$. These operations give
\[
\left[
\begin{array}{ccc@{\;\vline\;}ccc}
2&1&0&1&0&0\\
0&3&2&-1&2&0\\
0&0&4&1&-2&3
\end{array}
\right].
\]
Apply $R_3\leftarrow\tfrac14R_3$, then $R_2\leftarrow R_2-2R_3$, and finally $R_2\leftarrow\tfrac13R_2$:
\[
\left[
\begin{array}{ccc@{\;\vline\;}ccc}
2&1&0&1&0&0\\
0&1&0&-\tfrac12&1&-\tfrac12\\
0&0&1&\tfrac14&-\tfrac12&\tfrac34
\end{array}
\right].
\]
Finally, apply $R_1\leftarrow R_1-R_2$ and then divide row 1 by $2$:
\[
\left[
\begin{array}{ccc@{\;\vline\;}ccc}
1&0&0&\tfrac34&-\tfrac12&\tfrac14\\
0&1&0&-\tfrac12&1&-\tfrac12\\
0&0&1&\tfrac14&-\tfrac12&\tfrac34
\end{array}
\right].
\]
Therefore
\[
A^{-1}=\frac14
\begin{bmatrix}
3&-2&1\\
-2&4&-2\\
1&-2&3
\end{bmatrix}.
\]
For verification,
\[
\begin{bmatrix}
2&1&0\\
1&2&1\\
0&1&2
\end{bmatrix}
\begin{bmatrix}
3&-2&1\\
-2&4&-2\\
1&-2&3
\end{bmatrix}
=4I,
\]
so $AA^{-1}=I$.

\solutionquestion{7}
\solutionpart{i}Subtract row 1 from rows 2 and 3. The first-column multipliers are
\[
\ell_{21}=1,\qquad \ell_{31}=1.
\]
This gives
\[
\begin{bmatrix}
1&1&1\\
0&1&2\\
0&2&5
\end{bmatrix}.
\]
Next subtract twice row 2 from row 3, so $\ell_{32}=2$. The result is
\[
U=\begin{bmatrix}
1&1&1\\
0&1&2\\
0&0&1
\end{bmatrix}.
\]
Placing the multipliers in their corresponding positions gives
\[
L=\begin{bmatrix}
1&0&0\\
1&1&0\\
1&2&1
\end{bmatrix}.
\]
Direct multiplication confirms that
\[
LU=
\begin{bmatrix}
1&1&1\\
1&2&3\\
1&3&6
\end{bmatrix}=A.
\]

\solutionpart{ii}The system $Lc=b$ is
\[
\begin{aligned}
c_1&=5,\\
c_1+c_2&=7,\\
c_1+2c_2+c_3&=11.
\end{aligned}
\]
Forward substitution gives
\[
c=\begin{bmatrix}5\\2\\2\end{bmatrix}.
\]
The system $Ux=c$ is
\[
x_1+x_2+x_3=5,\qquad x_2+2x_3=2,\qquad x_3=2.
\]
Back substitution gives
\[
x_3=2,\qquad x_2=-2,\qquad x_1=5.
\]
Thus
\[
x=\begin{bmatrix}5\\-2\\2\end{bmatrix}.
\]

\solutionquestion{8}
\solutionpart{i}For $1\leq i\leq p$ and $1\leq j\leq m$,
\[
{\bigl({(AB)}^{T}\bigr)}_{ij}={(AB)}_{ji}
=\sum_{k=1}^{n}a_{jk}b_{ki}.
\]
On the other hand,
\[
\begin{aligned}
{(B^{T}A^{T})}_{ij}
&=\sum_{k=1}^{n}{(B^{T})}_{ik}{(A^{T})}_{kj}\\
&=\sum_{k=1}^{n}b_{ki}a_{jk}.
\end{aligned}
\]
The two sums are equal because real-number multiplication is commutative. Hence every entry agrees, and
\[
{(AB)}^{T}=B^{T}A^{T}.
\]

\solutionpart{ii}The $(i,j)$ entry of $Q^{T}Q$ is
\[
{(Q^{T}Q)}_{ij}=q_i^{T}q_j.
\]
Since $Q^{T}Q=I$,
\[
q_i^{T}q_j=\begin{cases}
1,&i=j,\\
0,&i\ne j.
\end{cases}
\]
Therefore each column has length $1$, and distinct columns are orthogonal. Conversely, if the columns are unit vectors and are mutually orthogonal, the same formula shows that every entry of $Q^{T}Q$ agrees with the corresponding entry of $I$. Thus $Q^{T}Q=I$.

\solutionpart{iii}For all $x,y\in\mathbb R^n$,
\[
{(Qx)}^{T}(Qy)=x^{T}Q^{T}Qy=x^{T}y.
\]
Thus $Q$ preserves dot products. Taking $y=x$ gives
\[
\|Qx\|^2={(Qx)}^{T}(Qx)=x^{T}x=\|x\|^2,
\]
so $Q$ also preserves lengths. A suitable $2\times2$ example is the rotation matrix
\[
Q=\begin{bmatrix}
\cos\theta&-\sin\theta\\
\sin\theta&\cos\theta
\end{bmatrix}.
\]
Its columns have unit length and zero dot product, so $Q^{T}Q=I$ and $q_{11}=\cos\theta$.

\end{document}
